\scanpage{16}\mspage{14}
2. \uline{Comparaison avec} \uline{Quillen}.

\deleted{\unsure{On}} Il se pose un certain nombre de questions sur ce foncteur $(M\Sigma^{-1})^o\times M\Sigma^{-1}\longrightarrow\mathrm{Hot}$, et sa relation aux constructions classiques, notamment à la théorie de Quillen. Je vais maintenant noter
\[
\begin{aligned}
\mathrm{Ch}_{\infty}^{\Sigma}(a,b)&=\mathrm{Hom}\,\mathrm{ext}_{\mathrm{Hot}}(a,b)\quad\text{ou simplement}\\
&\quad\mathrm{Hom}\,\mathrm{ext}(a,b).
\end{aligned}
\]

$1^{\mathrm o})$ Soient $f,g:a\rightrightarrows b$ dans $M$, alors (si $M$ est une ``model category'') Quillen définit un ensemble
\[
\pi_1(a,b;f,g)\quad\text{ou}\quad\pi_1(f\unsure{,|.}g)
\]
des classes d'homotopie de $f$ à \unsure{$g$}, de telle façon que $\mathrm{Hom}_{M}(a,b)$ devient \inserted{l'\unsure{ens.} des obj. d'} un groupoïde, dont les flèches sont les éléments des $\pi_1(f,g)$, pour une notion de composition naturelle.

Je suspecte que l'on a
\[
\pi_1(f,g)=\mathrm{Hom}_{\underbrace{\Pi(\mathrm{C}\underline{\mathrm{h}}_{\infty}^{\Sigma}(a,b))}_{\substack{\text{groupoïde}\\\text{fondamental}\\\text{de }\mathrm{C}\underline{\mathrm{h}}_{\infty}^{\Sigma}(a,b)}}}(f,g),
\]
quand $f$ et $g$ sont considérés comme des ``chemins infinis'' en \unsure{inscrivant|insérant} chaque $f:a\deleted{\unsure{\leftarrow}}\unsure{\text{---}|\to}b$ en
\diagram{\begin{tikzcd}[ampersand replacement=\&,column sep=small]
\ldots \& \underset{\substack{\Vert\\a}}{a_{-2}} \& \underset{\substack{\Vert\\a}}{a_{-1}} \arrow[l,"{\mathrm{id}}"'] \arrow[r,"f"] \& \underset{\substack{\Vert\\b}}{a_0} \& \underset{\substack{\Vert\\b}}{a_1} \arrow[l,"{\mathrm{id}}"'] \arrow[r] \& \ldots
\end{tikzcd}}
% archive sheet number: 15

